\section{Results}
\label{ch:results}

All results use \Rs10{,}000 invested on the first of every month and a 0.1\% cost on every trade.

\subsection{Design period: Feb 2010 -- Dec 2019}
\label{sec:res-2010}

Over 119 months, \Rs11.9 lakh was invested in each strategy.

\begin{table}[H]
\centering
\caption{Results, Feb 2010 -- Dec 2019 (in sample)}
\label{tab:res-2010}
\small
\begin{adjustbox}{max width=\textwidth}
\begin{tabular}{@{}lrrrrrr@{}}
\toprule
Strategy & Final value & XIRR & Volatility & Sharpe (vs 0\%) & Sharpe vs liquid & Max drawdown \\
\midrule
1 Equity & \Rs21.76 lakh & 11.7\% & 15.3\% & 0.76 & 0.29 & $-23.7\%$ \\
2 Equal-weight & \Rs18.50 lakh & 8.6\% & 6.5\% & 1.40 & 0.28 & $-6.5\%$ \\
3 60/20/20 & \Rs19.91 lakh & 10.0\% & 8.8\% & 1.15 & \textbf{0.33} & $-8.9\%$ \\
4 Risk parity & \Rs17.78 lakh & 7.9\% & 3.1\% & 2.59 & 0.26 & $-1.7\%$ \\
5 Max Sharpe & \Rs17.70 lakh & 7.8\% & 3.2\% & 2.49 & 0.24 & $-1.8\%$ \\
6 Optimized & \Rs17.74 lakh & 7.8\% & 3.2\% & 2.54 & 0.25 & $-1.7\%$ \\
\bottomrule
\end{tabular}
\end{adjustbox}
\end{table}

A SIP held entirely in the liquid fund earned 7.3\%.

\paragraph{Findings.}
\begin{enumerate}
  \item \textbf{100\% equity earned the most} (11.7\% a year) but suffered the largest fall
        ($-23.7\%$ in 2011).
  \item \textbf{60/20/20 gave the best return above the liquid fund per unit of risk} (0.33) and the
        second-highest return.
  \item \textbf{The optimisers (4--6) became about 70\% cash.} The liquid fund's volatility is
        about 0.5\%, so against a 0\% rate it appears to offer a high return for almost no risk. Both
        risk parity and maximum Sharpe therefore pushed it to the 70\% cap. Their Sharpe ratios
        against 0\% look excellent (about 2.5), but against the liquid fund they are the lowest of
        the six (0.24--0.26): most of their return is simply the T-bill rate. This is the issue
        \citet{sharpe1994} warns about, and it motivated Strategy 6b.
\end{enumerate}

\paragraph{Stress periods and robustness.} In four Indian stress periods (2011, the 2013 taper tantrum, the 2015--16 sell-off and the 2018 IL\&FS crisis), 100\% equity lost 8--24\%, while the gold- and liquid-heavy strategies were positive every time: rupee gold rose when the rupee fell \citep{baur2010}. Every 5-year SIP starting between 2010 and 2015 made money for every strategy.

\subsection{Test period: Jan 2020 -- Dec 2025}
\label{sec:res-2020}

Every rule was run unchanged on six years of data that played no part in its design (72
instalments, \Rs7.2 lakh). The winning rule, the highest Sharpe ratio versus liquid, was fixed before
the run.

\begin{table}[H]
\centering
\caption{Results, Jan 2020 -- Dec 2025 (out of sample)}
\label{tab:res-2020}
\small
\begin{adjustbox}{max width=\textwidth}
\begin{tabular}{@{}lrrrrr@{}}
\toprule
Strategy & Final value & XIRR & Volatility & Sharpe vs liquid & Max drawdown \\
\midrule
1 Equity & \Rs11.39 lakh & 15.1\% & 17.8\% & 0.55 & $-28.2\%$ \\
2 Equal-weight & \Rs12.02 lakh & 16.9\% & 7.3\% & \textbf{1.18} & $-6.8\%$ \\
3 60/20/20 & \Rs11.69 lakh & 16.0\% & 10.7\% & 0.83 & $-15.4\%$ \\
4 Risk parity & \Rs10.11 lakh & 11.2\% & 3.4\% & 1.16 & $-3.0\%$ \\
5 Max Sharpe & \Rs10.08 lakh & 11.1\% & 3.4\% & 1.13 & $-3.2\%$ \\
6 Optimized & \Rs10.09 lakh & 11.1\% & 3.4\% & 1.14 & $-3.1\%$ \\
6b Fixed Optimized & \textbf{\Rs13.28 lakh} & \textbf{20.3\%} & 10.3\% & 1.09 & $-12.2\%$ \\
\bottomrule
\end{tabular}
\end{adjustbox}
\end{table}

A SIP held entirely in the liquid fund earned 6.2\%.

\begin{figure}[H]
\centering
\includegraphics[width=0.9\textwidth]{08_forward_test/results/wealth_2020_2025.png}
\caption{Portfolio value of all strategies, Jan 2020 -- Dec 2025 (out of sample)}
\label{fig:wealth-2020}
\end{figure}

\paragraph{Findings.}
\begin{enumerate}
  \item \textbf{The equal-weight SIP won by the pre-set rule}, with the best return per unit of
        risk, consistent with \citet{demiguel2009}.
  \item \textbf{The correction in 6b worked out of sample.} The original Optimized SIP stayed about
        70\% in cash and earned 11.1\%; the fixed version earned 20.3\%, the highest of all.
  \item \textbf{Caveat: 2020--2025 was a gold boom.} Rupee gold rose about 21\% a year, including
        about 72\% in 2025. Gold-heavy mixes (equal-weight at one third, 6b at 33--49\%) gained the
        most. In 2010--2019, when rupee gold rose only about 3\% a year, neither of them was best.
\end{enumerate}

\paragraph{Cross-checks.} Rebuilt from the alternative funds of Section~\ref{ch:data}, every
strategy kept its rank and its XIRR changed by at most 0.17 percentage points. Strategies 1--3
recomputed with independent code matched the engine's final values exactly.

\subsection{Which approach maximises utility?}
\label{sec:res-utility}

Following the guide's suggestion, all strategies were then compared on one score that balances
return against risk, the utility $U$ of Eq.~\eqref{eq:U}, for three investor profiles. Strategy 7,
which maximises $U$ directly, was added to the competition.

\begin{table}[H]
\centering
\caption{Strategy with the highest utility $U$ for each investor profile}
\label{tab:utility}
\begin{adjustbox}{max width=\textwidth}
\begin{tabular}{@{}lll@{}}
\toprule
Investor & Highest $U$, 2010--2019 (in sample) & Highest $U$, 2020--2025 (out of sample) \\
\midrule
Aggressive ($\lambda=2$) & 60/20/20 (2.13\%) & Fixed Optimized (10.25\%) \\
Moderate ($\lambda=4$) & 60/20/20 (1.36\%) & Fixed Optimized (9.18\%) \\
Conservative ($\lambda=8$) & Risk parity (0.44\%) & Fixed Optimized (7.04\%) \\
\bottomrule
\end{tabular}
\end{adjustbox}
\end{table}

Strategy 7 ranked 4th, 4th and 6th in 2010--2019 and 2nd, 2nd and 3rd in 2020--2025 (for $\lambda = 2, 4, 8$).

\paragraph{Findings.}
\begin{enumerate}
  \item \textbf{The best strategy depends on the investor.} In 2010--2019 the answer changes with
        $\lambda$.
  \item \textbf{Maximising $U$ directly did not maximise the $U$ achieved.} Strategy 7 optimises
        exactly the score it is judged on, yet it never came first. It trusts the window's average
        returns, which are noisy, and over-weights whatever did well recently. This is the
        estimation-error problem described by \citet{michaud1989} and \citet{chopra1993}. Rules that
        rely less on average returns held up better.
  \item \textbf{The Fixed Optimized SIP had the highest $U$ for every profile out of sample}, and
        the ranking held on the alternative-fund data. In 2010--2019 it ranked 3rd, 2nd and 7th.
  \item \textbf{100\% equity never won}: its 17.8\% volatility in 2020--2025 cost more, in $U$,
        than its extra return.
\end{enumerate}

\subsection{The Personal SIP}
\label{sec:res-personal}

\subsubsection{Choosing the cushion (2010--2019 only)}

\begin{table}[H]
\centering
\caption{$U$ of the Personal SIP for each liquid cushion, Feb 2010 -- Dec 2019 (best in bold)}
\label{tab:cushion}
\begin{adjustbox}{max width=\textwidth}
\begin{tabular}{@{}cccc@{}}
\toprule
Cushion & Aggressive ($\lambda=2$) & Moderate ($\lambda=4$) & Conservative ($\lambda=8$) \\
\midrule
10\% & \textbf{1.52\%} & 0.68\% & $-0.99\%$ \\
20\% & 1.41\% & 0.75\% & $-0.57\%$ \\
30\% & 1.31\% & \textbf{0.80\%} & $-0.21\%$ \\
40\% & 1.17\% & 0.80\% & 0.05\% \\
50\% & 1.01\% & 0.75\% & 0.23\% \\
60\% & 0.83\% & 0.66\% & 0.33\% \\
70\% & 0.69\% & 0.59\% & \textbf{0.40\%} \\
\bottomrule
\end{tabular}
\end{adjustbox}
\end{table}

The chosen cushions are 10\% (aggressive), 30\% (moderate) and 70\% (conservative). For the
moderate investor, 30\% narrowly beat 40\% (0.7997\% against 0.7990\%). The calculation for the
moderate investor shows the trade-off that $U$ resolves:
\begin{align*}
  L = 10\%:&\quad U = 2.354\% - 2 \times (9.145\%)^2 = 2.354\% - 1.673\% = 0.681\% \\
  L = 30\%:&\quad U = 1.814\% - 2 \times (7.121\%)^2 = 1.814\% - 1.014\% = 0.800\% \\
  L = 70\%:&\quad U = 0.786\% - 2 \times (3.096\%)^2 = 0.786\% - 0.192\% = 0.594\%
\end{align*}
A small cushion earns the most but swings so much that the penalty absorbs most of the gain; a large
cushion barely swings but barely earns.

\subsubsection{Out-of-sample test (2020--2025)}

\begin{table}[H]
\centering
\caption{Personal SIP versus the Fixed Optimized SIP, Jan 2020 -- Dec 2025}
\label{tab:personal}
\small
\begin{adjustbox}{max width=\textwidth}
\begin{tabular}{@{}lcccc@{}}
\toprule
Investor (cushion) & \multicolumn{2}{c}{Personal SIP} & Rank & Fixed Optimized SIP \\
 & $U$ & XIRR / worst fall & (of 11) & $U$ / XIRR / worst fall \\
\midrule
Aggressive (10\%) & 9.99\% & 20.4\% / $-11.8\%$ & 3 & 10.25\% / 20.3\% / $-12.2\%$ \\
Moderate (30\%) & 7.25\% & 17.2\% / $-9.0\%$ & 5 & 9.18\% / 20.3\% / $-12.2\%$ \\
Conservative (70\%) & 3.24\% & 11.0\% / $-3.4\%$ & 9 & 7.04\% / 20.3\% / $-12.2\%$ \\
\bottomrule
\end{tabular}
\end{adjustbox}
\end{table}

\paragraph{Findings.}
\begin{enumerate}
  \item \textbf{It personalises the advice.} The conservative investor's worst fall shrinks from
        $-12.2\%$ to $-3.4\%$.
  \item \textbf{It did not beat the Fixed Optimized SIP on $U$.} The cushions were chosen in
        2010--2019, when Nifty and gold earned little above the liquid fund; 2020--2025 rewarded
        holding them. A cushion learned in one period can be too cautious for the next, the same
        estimation problem as in Section~\ref{sec:res-utility}. Even for $\lambda = 8$, the Fixed
        Optimized SIP's extra 7.6\% a year above liquid outweighed its larger variance penalty
        ($4 \times 10.34\%^2 = 4.28\%$ against $4 \times 3.38\%^2 = 0.46\%$).
  \item \textbf{Re-fitting monthly or yearly gave almost the same results}, so the monthly rolling
        window used by the website costs nothing in performance.
\end{enumerate}

The Personal SIP was nevertheless chosen for the website because it is the only rule in which the
investor's stated risk tolerance changes the advice, which is the purpose of asking for it.

\paragraph{Live check, Jan--Oct 2026.} In ten months published after all choices were made, Nifty fell (March $-11.3\%$) and gold swung; the cash-heavy Optimized SIP ended slightly ahead (\Rs1{,}00{,}546 on \Rs1{,}00{,}000), the gold-heavy Fixed Optimized SIP slightly behind (\Rs95{,}748). Ten months is too short to judge a SIP.

